\documentclass[10pt,a4paper]{amsart}
\usepackage{amsmath,amssymb,amsthm,graphicx}
% Retained minimal preamble for the Stacks Project chapter varieties.tex.
% Verbatim excerpt of the chapter preamble: Xy-pic setup (native preamble lines 14-18), the native theorem
% declarations with the mathematical macros (native preamble lines 154-217), and the project own no-op
% definitions of the reference/slogan/history display environments (native wrapper.tex lines 1-3) that the
% counted statement uses for its slogan. Copied mechanically; nothing added, removed or redefined.
% For commutative diagrams we use Xy-pic
\usepackage[all]{xy}

% We use 2cell for 2-commutative diagrams.
\xyoption{2cell}
% Theorem environments.
%
\theoremstyle{plain}
\newtheorem{theorem}[subsection]{Theorem}
\newtheorem{proposition}[subsection]{Proposition}
\newtheorem{lemma}[subsection]{Lemma}

\theoremstyle{definition}
\newtheorem{definition}[subsection]{Definition}
\newtheorem{example}[subsection]{Example}
\newtheorem{exercise}[subsection]{Exercise}
\newtheorem{situation}[subsection]{Situation}

\theoremstyle{remark}
\newtheorem{remark}[subsection]{Remark}
\newtheorem{remarks}[subsection]{Remarks}

\numberwithin{equation}{subsection}

% Macros
%
\def\lim{\mathop{\mathrm{lim}}\nolimits}
\def\colim{\mathop{\mathrm{colim}}\nolimits}
\def\Spec{\mathop{\mathrm{Spec}}}
\def\Hom{\mathop{\mathrm{Hom}}\nolimits}
\def\Ext{\mathop{\mathrm{Ext}}\nolimits}
\def\SheafHom{\mathop{\mathcal{H}\!\mathit{om}}\nolimits}
\def\SheafExt{\mathop{\mathcal{E}\!\mathit{xt}}\nolimits}
\def\Sch{\mathit{Sch}}
\def\Mor{\mathop{\mathrm{Mor}}\nolimits}
\def\Ob{\mathop{\mathrm{Ob}}\nolimits}
\def\Sh{\mathop{\mathit{Sh}}\nolimits}
\def\NL{\mathop{N\!L}\nolimits}
\def\CH{\mathop{\mathrm{CH}}\nolimits}
\def\proetale{{pro\text{-}\acute{e}tale}}
\def\etale{{\acute{e}tale}}
\def\QCoh{\mathit{QCoh}}
\def\Ker{\mathop{\mathrm{Ker}}}
\def\Im{\mathop{\mathrm{Im}}}
\def\Coker{\mathop{\mathrm{Coker}}}
\def\Coim{\mathop{\mathrm{Coim}}}

% Boxtimes
%
\DeclareMathSymbol{\boxtimes}{\mathbin}{AMSa}{"02}

%
% Macros for moduli stacks/spaces
%
\def\QCohstack{\mathcal{QC}\!\mathit{oh}}
\def\Cohstack{\mathcal{C}\!\mathit{oh}}
\def\Spacesstack{\mathcal{S}\!\mathit{paces}}
\def\Quotfunctor{\mathrm{Quot}}
\def\Hilbfunctor{\mathrm{Hilb}}
\def\Curvesstack{\mathcal{C}\!\mathit{urves}}
\def\Polarizedstack{\mathcal{P}\!\mathit{olarized}}
\def\Complexesstack{\mathcal{C}\!\mathit{omplexes}}
% \Pic is the operator that assigns to X its picard group, usage \Pic(X)
% \Picardstack_{X/B} denotes the Picard stack of X over B
% \Picardfunctor_{X/B} denotes the Picard functor of X over B
\def\Pic{\mathop{\mathrm{Pic}}\nolimits}
\def\Picardstack{\mathcal{P}\!\mathit{ic}}
\def\Picardfunctor{\mathrm{Pic}}
\def\Deformationcategory{\mathcal{D}\!\mathit{ef}}
\newenvironment{reference}{}{}
\newenvironment{slogan}{}{}
\newenvironment{history}{}{}
\usepackage{cleveref}
\crefname{originaltheorem}{original theorem}{original theorems}
\Crefname{originaltheorem}{Original theorem}{Original theorems}
\makeatletter
\newlabel{algebra-theorem-nullstellensatz}{{00FV}{source}{}{sourcepointer-1}{}}
\newlabel{algebra-theorem-nullstellensatz@cref}{{[originaltheorem][0][]00FV}{[1][0][]source}{}{}{}}
\makeatother
\crefname{originallemma}{original lemma}{original lemmas}
\Crefname{originallemma}{Original lemma}{Original lemmas}
\makeatletter
\newlabel{divisors-lemma-push-down-quasi-affine}{{0BD5}{source}{}{sourcepointer-2}{}}
\newlabel{divisors-lemma-push-down-quasi-affine@cref}{{[originallemma][0][]0BD5}{[1][0][]source}{}{}{}}
\makeatother
\crefname{originallemma}{original lemma}{original lemmas}
\Crefname{originallemma}{Original lemma}{Original lemmas}
\makeatletter
\newlabel{limits-lemma-limit-quasi-affine}{{01Z5}{source}{}{sourcepointer-3}{}}
\newlabel{limits-lemma-limit-quasi-affine@cref}{{[originallemma][0][]01Z5}{[1][0][]source}{}{}{}}
\makeatother
\crefname{originaldefinition}{original definition}{original definitions}
\Crefname{originaldefinition}{Original definition}{Original definitions}
\makeatletter
\newlabel{properties-definition-quasi-affine}{{01P6}{source}{}{sourcepointer-4}{}}
\newlabel{properties-definition-quasi-affine@cref}{{[originaldefinition][0][]01P6}{[1][0][]source}{}{}{}}
\makeatother
\crefname{originallemma}{original lemma}{original lemmas}
\Crefname{originallemma}{Original lemma}{Original lemmas}
\makeatletter
\newlabel{properties-lemma-quasi-affine-locally-closed}{{0BCK}{source}{}{sourcepointer-5}{}}
\newlabel{properties-lemma-quasi-affine-locally-closed@cref}{{[originallemma][0][]0BCK}{[1][0][]source}{}{}{}}
\makeatother
\crefname{originallemma}{original lemma}{original lemmas}
\Crefname{originallemma}{Original lemma}{Original lemmas}
\makeatletter
\newlabel{schemes-lemma-characterize-quasi-separated}{{01KO}{source}{}{sourcepointer-6}{}}
\newlabel{schemes-lemma-characterize-quasi-separated@cref}{{[originallemma][0][]01KO}{[1][0][]source}{}{}{}}
\makeatother
\title{Quasi-affineness descends along an arbitrary field extension}
\author{The Stacks Project Authors}
\date{stacks-project:varieties.tex:local-snapshot-20261011; GNU FDL 1.2 or later}
\begin{document}
\maketitle
\noindent\textit{Editorial note.} All delivered mathematical text is the author's own, copied from the pinned Stacks Project chapter file varieties.tex; no statement or proof step was rewritten, reordered or removed.\par\medskip
\noindent\textit{Editorial note.} The proof is complete as written and ends with the descent along the finite locally free morphism; no step is referred to the reader.\par\medskip
\noindent\textit{Editorial note.} The counted spans carry 6 navigation stubs and no bibliography citation.\par\medskip
\noindent\textit{Editorial note.} The preamble is the chapter's retained minimal preamble (sources/stacks-varieties/retained-preamble.tex lines 1-77), a verbatim excerpt of the native preamble.\par\medskip
\noindent\textit{Editorial note.} The label lemma-quasi-affine-after-field-extension is defined only in varieties.tex, so no cross-chapter restatement had to be weighed for this card.\par\medskip
\noindent\textit{Allowed standard prerequisites (editorial).}\begin{itemize}
\item Quasi-affine means quasi-compact and an open subscheme of an affine scheme. Properties, Definition properties-definition-quasi-affine, used at proof lines 2790-2791 to extract quasi-compactness. Supplies that X\_K, being quasi-affine, is quasi-compact, so that its image X is quasi-compact.
\item Criterion for quasi-separatedness through affine opens. Schemes, Lemma schemes-lemma-characterize-quasi-separated, invoked at proof lines 2796-2797. Supplies quasi-separatedness of X.
\item Quasi-affineness descends through a filtered limit. Limits, Lemma limits-lemma-limit-quasi-affine, invoked at proof lines 2803-2804. Supplies the index i for which X\_i is already quasi-affine.
\item Quasi-affineness is local for the Zariski topology and descends along a finite locally free morphism. Properties, Lemma properties-lemma-quasi-affine-locally-closed and Divisors, Lemma divisors-lemma-push-down-quasi-affine, invoked at proof lines 2809-2813. Supplies that X\_\{k'\} is quasi-affine as a closed subscheme and that quasi-affineness descends from X\_\{k'\} to X.
\end{itemize}\medskip
\ifdefined\hypertarget\hypertarget{sourcepointer-1}{}\fi
\noindent\textit{Source reference.} Original theorem 00FV: sources/stacks-chow/tags.txt, source lines 585--585. Pointer outside the counted span, invoked once for context; the load-bearing argument of the counted statement is the author's proof reproduced in full in the proof span.\par\medskip
\ifdefined\hypertarget\hypertarget{sourcepointer-2}{}\fi
\noindent\textit{Source reference.} Original lemma 0BD5: sources/stacks-chow/tags.txt, source lines 14303--14303. Pointer outside the counted span, invoked once for context; the load-bearing argument of the counted statement is the author's proof reproduced in full in the proof span.\par\medskip
\ifdefined\hypertarget\hypertarget{sourcepointer-3}{}\fi
\noindent\textit{Source reference.} Original lemma 01Z5: sources/stacks-chow/tags.txt, source lines 2580--2580. Pointer outside the counted span, invoked once for context; the load-bearing argument of the counted statement is the author's proof reproduced in full in the proof span.\par\medskip
\ifdefined\hypertarget\hypertarget{sourcepointer-4}{}\fi
\noindent\textit{Source reference.} Original definition 01P6: sources/stacks-chow/tags.txt, source lines 2219--2219. Pointer outside the counted span, invoked once for context; the load-bearing argument of the counted statement is the author's proof reproduced in full in the proof span.\par\medskip
\ifdefined\hypertarget\hypertarget{sourcepointer-5}{}\fi
\noindent\textit{Source reference.} Original lemma 0BCK: sources/stacks-chow/tags.txt, source lines 14283--14283. Pointer outside the counted span, invoked once for context; the load-bearing argument of the counted statement is the author's proof reproduced in full in the proof span.\par\medskip
\ifdefined\hypertarget\hypertarget{sourcepointer-6}{}\fi
\noindent\textit{Source reference.} Original lemma 01KO: sources/stacks-chow/tags.txt, source lines 2057--2057. Pointer outside the counted span, invoked once for context; the load-bearing argument of the counted statement is the author's proof reproduced in full in the proof span.\par\medskip
\begin{lemma}
\label{lemma-quasi-affine-after-field-extension}
Let $k$ be a field. Let $X$ be a scheme over $k$. If $X_K$ is quasi-affine
for some field extension $K/k$, then $X$ is quasi-affine.
\end{lemma}

\begin{proof}
Let $K/k$ be a field extension such that $X_K$ is quasi-affine.
The morphism $X_K \to X$ is surjective. Hence $X$ is quasi-compact
as the image of a quasi-compact scheme (Properties, Definition
\ref{properties-definition-quasi-affine}). Since $X_K$ is quasi-separated
(as an open subscheme of an affine scheme)
we see that $X$ is quasi-separated: If $U, V \subset X$ are
affine open, then $(U \cap V)_K = U_K \cap V_K$ is quasi-compact
and $(U \cap V)_K \to U \cap V$ is surjective. Thus
Schemes, Lemma \ref{schemes-lemma-characterize-quasi-separated} applies.

\medskip\noindent
Write $K = \colim A_i$ as the colimit of the subalgebras of $K$
which are of finite type over $k$. Denote
$X_i = X \times_{\Spec(k)} \Spec(A_i)$.
Since $X_K = \lim X_i$ we find an $i$ such that $X_i$ is quasi-affine
(Limits, Lemma \ref{limits-lemma-limit-quasi-affine};
here we use that $X$ is quasi-compact and quasi-separated as
just shown). By the Hilbert Nullstellensatz
(Algebra, Theorem \ref{algebra-theorem-nullstellensatz})
the residue field $k' = A_i/\mathfrak m$ is a finite
extension of $k$. Hence $X_{k'} \subset X_i$ is a closed
subscheme hence is quasi-affine (Properties, Lemma
\ref{properties-lemma-quasi-affine-locally-closed}).
Since $X_{k'} \to X$ is finite locally free we conclude by
Divisors, Lemma \ref{divisors-lemma-push-down-quasi-affine}.
\end{proof}

\end{document}
